\section{Theoretical Analysis}
\label{sec:theory}

We analyze one head and suppress batch and head indices. Let
\begin{equation}
  A=\mu\I_N+\gamma UU^\top,
  \qquad \mu>0,\quad\gamma\geq0.
\end{equation}
The same statements apply to \RRLSSO\ after replacing $U$ by $\widetilde U$.

\subsection{Well-Posedness and Exact Rank-Space Solution}

\begin{proposition}[Existence and uniqueness]
\label{prop:spd}
$A$ is symmetric positive definite. Hence $AY=C$ has a unique solution for every $C$.
\end{proposition}
\begin{proof}
For every nonzero $z\in\R^N$,
\[
  z^\top Az
  =\mu\lVert z\rVert_2^2
  +\gamma\lVert U^\top z\rVert_2^2
  \geq\mu\lVert z\rVert_2^2>0.
\]
\end{proof}

\begin{proposition}[Exact Woodbury reduction]
\label{prop:woodbury}
The unique solution is given by \cref{eq:woodbury}, including at $\gamma=0$.
\end{proposition}
\begin{proof}
For $\gamma>0$, apply the Woodbury identity to $\mu\I_N+U(\gamma\I_r)U^\top$:
\[
  A^{-1}
  =\mu^{-1}\I_N
  -\mu^{-1}U
  (\gamma^{-1}\I_r+\mu^{-1}U^\top U)^{-1}
  U^\top\mu^{-1}.
\]
Multiplying the inner inverse by $\gamma$ gives \cref{eq:woodbury}. The resulting expression is continuous at $\gamma=0$ and then reduces to $Y=\mu^{-1}C$.
\end{proof}

The proposition separates two notions that are sometimes conflated. The induced correction is low-rank, but the rank-space computation is not an approximation: it evaluates the full inverse operator exactly.

\subsection{Trace-Normalized Relation Basis}

Let $Z\in\R^{N\times r}$ be the raw relation projection and set $\tau=rL_{\mathrm{ref}}$. We define
\begin{equation}
  U=\mathcal N_{\mathrm{tr}}(Z)
  =\frac{\sqrt{\tau}}{\lVert Z\rVert_F}Z,
  \label{eq:trace_normalization}
\end{equation}
with masked entries excluded from the norm and set to zero. The implementation adds the usual RMS-style epsilon to the squared norm for the degenerate neighborhood of $Z=0$.

\begin{proposition}[Minimal scale fixing]
Among positive scalar maps $\widehat Z=a(Z)Z$ that preserve every token-to-token norm ratio and satisfy $\lVert\widehat Z\rVert_F^2=\tau$, \cref{eq:trace_normalization} is unique. Moreover,
\begin{equation}
  \operatorname{tr}(UU^\top)=\tau,
  \qquad
  \lambda_{\max}(UU^\top)\leq\tau.
\end{equation}
\end{proposition}
\begin{proof}
The trace constraint gives $a(Z)^2\lVert Z\rVert_F^2=\tau$, whose unique positive solution is $a(Z)=\sqrt{\tau}/\lVert Z\rVert_F$. The trace identity follows from $\operatorname{tr}(UU^\top)=\lVert U\rVert_F^2$, and the spectral bound follows because a positive-semidefinite matrix's largest eigenvalue does not exceed its trace.
\end{proof}

Unlike row-wise RMS normalization, this removes only the single global scale that is non-identifiable with $\gamma$: replacing $Z$ by $aZ$ and $\gamma$ by $\gamma/a^2$ leaves the solve unchanged. It preserves the radial allocation
\[
  (UU^\top)_{ii}
  =\tau\frac{\lVert z_i\rVert_2^2}{\lVert Z\rVert_F^2},
\]
so token-dependent relation strength remains expressible. Finally, row-wise orthogonal rank rotations commute with $\mathcal N_{\mathrm{tr}}$ because they preserve $\lVert Z\rVert_F$.

\subsection{Spectral Response, Stability, and Signed Mixing}

Let $UU^\top=Q\Lambda Q^\top$ over its $k$ nonzero eigen-directions, where $Q^\top Q=\I_k$ and $\Lambda=\diag(\lambda_1,\ldots,\lambda_k)$ with $\lambda_i>0$. Define $P_\perp=\I_N-QQ^\top$.

\begin{proposition}[Input-conditioned inverse filter]
\label{prop:spectrum}
For fixed $U$,
\begin{equation}
  A^{-1}
  =Q\diag\!\left(\frac{1}{\mu+\gamma\lambda_i}\right)Q^\top
  +\mu^{-1}P_\perp.
  \label{eq:spectral_filter}
\end{equation}
Thus \LSSO\ applies selective shrinkage on the relation subspace and isotropic scale $\mu^{-1}$ on its orthogonal complement.
\end{proposition}
\begin{proof}
On $\operatorname{span}(Q)$, $A$ has eigenvalues $\mu+\gamma\lambda_i$. On its orthogonal complement, $UU^\top$ vanishes and $A$ acts as $\mu\I$.
\end{proof}

\begin{proposition}[Norm and conditioning bounds]
\label{prop:bounds}
For fixed $U$,
\begin{align}
  \lVert Y\rVert_F&\leq\mu^{-1}\lVert C\rVert_F,\label{eq:norm_bound}\\
  \kappa_2(A)&\leq1+\frac{\gamma}{\mu}\lVert U\rVert_2^2.
  \label{eq:condition_bound}
\end{align}
For multiple heads followed by $W_O$,
\begin{equation}
  \lVert Y_{\mathrm{cat}}W_O\rVert_F
  \leq\mu_{\min}^{-1}
  \lVert C_{\mathrm{cat}}\rVert_F\lVert W_O\rVert_2.
\end{equation}
\end{proposition}
\begin{proof}
Since $A\succeq\mu\I$, $\lVert A^{-1}\rVert_2\leq\mu^{-1}$, proving \cref{eq:norm_bound}. Moreover,
\[
  \lambda_{\min}(A)\geq\mu,
  \qquad
  \lambda_{\max}(A)\leq\mu+\gamma\lVert U\rVert_2^2,
\]
which gives \cref{eq:condition_bound}. The multi-head result follows by summing squared head-wise bounds and applying submultiplicativity to $W_O$.
\end{proof}

These are conditional bounds for the linear map at fixed $U$; because both $U$ and $C$ depend on $X$, they are not by themselves a global Lipschitz bound for the complete nonlinear layer. Trace normalization gives the explicit bound $\lVert U\rVert_2^2\leq rL_{\mathrm{ref}}$ without forcing every token to have the same radius.

Although $UU^\top$ and $A^{-1}$ are positive semidefinite and positive definite respectively, their off-diagonal entries need not be nonnegative. For example, with two tokens and rank one, $U=[1,1]^\top$ yields a negative off-diagonal entry in $A^{-1}$, whereas $U=[1,-1]^\top$ yields a positive one. Thus the effective mixing operator supports signed cross-token coefficients and is not constrained to be a row-stochastic nonnegative matrix. Positive definiteness constrains the global quadratic form, not the sign of every pairwise interaction.

\subsection{Strongly Convex Energy}

\begin{theorem}[Variational characterization]
\label{thm:energy}
The \LSSO\ output is the unique minimizer
\begin{equation}
  Y^\star
  =\arg\min_Y
  \left[
    \frac{\mu}{2}\lVert Y\rVert_F^2
    +\frac{\gamma}{2}\lVert U^\top Y\rVert_F^2
    -\langle C,Y\rangle
  \right].
  \label{eq:energy}
\end{equation}
Equivalently, up to a constant independent of $Y$,
\begin{equation}
  Y^\star
  =\arg\min_Y
  \left[
    \frac{\mu}{2}\left\lVert Y-\mu^{-1}C\right\rVert_F^2
    +\frac{\gamma}{2}\lVert U^\top Y\rVert_F^2
  \right].
  \label{eq:energy_centered}
\end{equation}
\end{theorem}
\begin{proof}
The Hessian with respect to each output channel is $A$, which is positive definite by \cref{prop:spd}; the objective is therefore strongly convex. Its first-order condition is
\[
  \mu Y+\gamma UU^\top Y-C=0,
\]
which is exactly \cref{eq:lsso}. Completing the square gives \cref{eq:energy_centered}.
\end{proof}

The content projection supplies the unconstrained state $\mu^{-1}C$. The second term imposes an input-conditioned low-rank penalty, and the solve returns the unique balance between content fidelity and relation-space suppression. Unlike the Neumann interpretation below, this result requires no contraction condition.

\subsection{Rank-Space Inner Learning and Primal--Dual Equivalence}

The Woodbury variable is the solution of a sample-conditioned learning
problem, not merely an algebraic intermediate.

\begin{theorem}[Closed-form inner learner]
\label{thm:inner_learner}
Let $\alpha=\gamma/\mu$. The objective
\begin{equation}
  \mathcal J(Z)
  =\frac12\lVert Z\rVert_F^2
  +\frac{\alpha}{2}\lVert UZ-C\rVert_F^2
  \label{eq:rank_ridge}
\end{equation}
is strongly convex and has the unique minimizer
\begin{equation}
  Z^\star
  =(\I_r+\alpha U^\top U)^{-1}\alpha U^\top C.
  \label{eq:z_star}
\end{equation}
The corresponding \LSSO\ output satisfies
\begin{equation}
  \mu Y^\star=C-UZ^\star,
  \qquad
  Z^\star=\gamma U^\top Y^\star.
  \label{eq:primal_dual_links}
\end{equation}
Consequently, the token-space energy in \cref{eq:energy} and rank-space ridge
problem in \cref{eq:rank_ridge} are an exact primal--dual pair.
\end{theorem}
\begin{proof}
The Hessian of $\mathcal J$ is
$\I_r+\alpha U^\top U\succ0$. Its first-order condition is
\[
  (\I_r+\alpha U^\top U)Z=\alpha U^\top C,
\]
which gives \cref{eq:z_star}. Substitution into the Woodbury formula gives the
first identity in \cref{eq:primal_dual_links}. Left-multiplying it by
$\gamma U^\top$ and using the normal equation gives the second identity.
For $\gamma>0$, equivalently introducing a Fenchel variable for
$(\gamma/2)\lVert U^\top Y\rVert_F^2$ in \cref{eq:energy} produces the dual
minimization
\[
  \min_Z
  \left[
    \frac{1}{2\gamma}\lVert Z\rVert_F^2
    +\frac{1}{2\mu}\lVert UZ-C\rVert_F^2
  \right],
\]
which differs from \cref{eq:rank_ridge} only by the positive factor $\gamma$.
At $\gamma=0$, both formulations reduce continuously to $Z^\star=0$ and
$Y^\star=\mu^{-1}C$.
\end{proof}

For each input, the outer network generates $U(X)$ and $C(X)$ and thereby
constructs a new ridge problem. Its optimum $Z^\star(X)$ is a matrix of
per-sample fast weights. The relation is exact for the shared linear learner in
\cref{eq:rank_ridge}; it is not an equivalence to arbitrary multilayer,
untied, or nonlinear inner networks.

\subsection{Regularized Residual Readout}

The fitted component of the inner learner is $\widehat C=UZ^\star$, whereas
\LSSO\ returns
\begin{equation}
  Y^\star=\mu^{-1}(C-\widehat C).
  \label{eq:ridge_residual}
\end{equation}
For $\alpha>0$, this is a ridge residual, not an orthogonal least-squares
residual: the normal equation implies
\begin{equation}
  U^\top(C-UZ^\star)=\alpha^{-1}Z^\star,
\end{equation}
which is generally nonzero. At $\alpha=0$, $Z^\star=0$ and the residual is
simply $C$. Nevertheless, \cref{eq:spectral_filter} gives an
exact interpretation. Modes with large relation singular values are strongly
suppressed, weakly represented modes are suppressed less, and components
orthogonal to the relation subspace retain scale $\mu^{-1}$. Thus the operator
removes sample-predictable shared structure through a regularized shrinkage
path while preserving weakly explained and orthogonal information.

\subsection{Inner Optimization Limit and Neumann Special Case}

Let
\begin{equation}
  H=\I_r+\alpha U^\top U,
  \qquad b=\alpha U^\top C.
\end{equation}
Gradient descent on \cref{eq:rank_ridge} is the Richardson iteration
\begin{equation}
  Z^{(t+1)}=Z^{(t)}-\eta(HZ^{(t)}-b).
  \label{eq:inner_iteration}
\end{equation}

\begin{proposition}[Exact infinite-step inner-training limit]
\label{prop:inner_limit}
For every initialization $Z^{(0)}$, if
\begin{equation}
  0<\eta<\frac{2}{\lambda_{\max}(H)},
  \label{eq:richardson_condition}
\end{equation}
then \cref{eq:inner_iteration} converges to $Z^\star$. Hence the structured
solve evaluates exactly the infinite-step limit of any such convergent inner
training process.
\end{proposition}
\begin{proof}
Subtracting $Z^\star=H^{-1}b$ from the iteration gives
\[
  Z^{(t+1)}-Z^\star
  =(\I_r-\eta H)(Z^{(t)}-Z^\star).
\]
Because $H\succ0$, condition \cref{eq:richardson_condition} implies
$\rho(\I_r-\eta H)<1$.
\end{proof}

This result removes the inner update count, learning-rate choice, activation
storage, and truncation error for the convex learner. It eliminates
\emph{optimization depth}; it does not make a deeper nonlinear inner
architecture analytically solvable.

When $\eta=1$ and $Z^{(0)}=0$, \cref{eq:inner_iteration} becomes
\begin{equation}
  Z^{(t+1)}
  =\alpha U^\top C-\alpha U^\top UZ^{(t)}.
  \label{eq:unit_inner_iteration}
\end{equation}

\begin{proposition}[Neumann as unit-step inner adaptation]
\label{prop:neumann}
If
\begin{equation}
  \alpha\lambda_{\max}(U^\top U)<1,
  \label{eq:neumann_condition}
\end{equation}
then the unit-step iteration converges and
\begin{align}
  Z^\star
  &=\alpha\sum_{k=0}^{\infty}
    (-\alpha U^\top U)^kU^\top C,
    \label{eq:rank_neumann}\\
  Y^\star
  &=\frac{1}{\mu}
    \sum_{k=0}^{\infty}(-\alpha UU^\top)^kC.
    \label{eq:neumann}
\end{align}
\end{proposition}
\begin{proof}
Condition \cref{eq:neumann_condition} makes
$-\alpha U^\top U$ contractive. Unrolling
\cref{eq:unit_inner_iteration} gives \cref{eq:rank_neumann}; substitution into
$Y^\star=\mu^{-1}(C-UZ^\star)$ yields \cref{eq:neumann}.
\end{proof}

Neumann is therefore one particular inner optimizer: zero initialization,
unit step size, and a stronger spectral condition. General Richardson
iteration admits any safe step in \cref{eq:richardson_condition}, while the
positive-definite solve bypasses iteration altogether and remains defined for
every $\mu>0$ and $\gamma\geq0$. A depth-$K$ Neumann truncation realizes a
polynomial response; \LSSO\ directly realizes the rational response
$g(\lambda)=(\mu+\gamma\lambda)^{-1}$.

\subsection{Guarantees for Rank-Rotary LSSO}

The preceding results guarantee that rank rotation does not destroy the
positive-definite solve. We next characterize what position-dependent
row-wise rotation preserves and what it can change. Let the $i$-th row be
represented as a column vector $u_i\in\R^r$, and define
\begin{equation}
  \widetilde u_i=R_i u_i,
  \qquad R_i^\top R_i=\I_r.
  \label{eq:rowwise_rotation}
\end{equation}

\begin{lemma}[Spectral-energy preservation under row-wise rotation]
\label{lem:rotation_energy}
For arbitrary token-dependent orthogonal matrices $R_i$,
\begin{align}
  \lVert\widetilde U\rVert_F^2
  &=\lVert U\rVert_F^2,\label{eq:frobenius_preserved}\\
  \tr(\widetilde U\widetilde U^\top)
  &=\tr(UU^\top).\label{eq:trace_preserved}
\end{align}
Consequently, row-wise rotation preserves the total relation spectral energy,
although it need not preserve the individual singular values.
\end{lemma}
\begin{proof}
Orthogonality preserves every row norm:
\[
  \lVert\widetilde u_i\rVert_2^2
  =u_i^\top R_i^\top R_i u_i
  =\lVert u_i\rVert_2^2.
\]
Summing over tokens gives \cref{eq:frobenius_preserved}. The trace identity
follows from
$\tr(\widetilde U\widetilde U^\top)=\lVert\widetilde U\rVert_F^2$
and the corresponding identity for $U$.
\end{proof}

The preservation is global rather than directional. For example, take two
identical unit rows in $\R^2$. Before rotation the singular values are
$(\sqrt{2},0)$. Rotating only the second row to the orthogonal coordinate gives
singular values $(1,1)$. Both bases have squared Frobenius norm $2$, but their
spectral concentration differs.

Define the spectral factor governing the rank-rotary fixed-point iteration by
\begin{equation}
  q_{\max}^{R}
  =\rho(\alpha\widetilde U\widetilde U^\top)
  =\frac{\gamma}{\mu}
   \lambda_{\max}(\widetilde U\widetilde U^\top)
  =\frac{\gamma}{\mu}\lVert\widetilde U\rVert_2^2.
  \label{eq:qmax_rotary}
\end{equation}
For a nonzero matrix, let its stable rank be
\begin{equation}
  r_s(\widetilde U)
  =\frac{\lVert\widetilde U\rVert_F^2}
  {\lVert\widetilde U\rVert_2^2}.
  \label{eq:stable_rank}
\end{equation}
Combining \cref{eq:frobenius_preserved,eq:qmax_rotary,eq:stable_rank} yields
\begin{equation}
  q_{\max}^{R}
  =\frac{\gamma}{\mu}
   \frac{\lVert U\rVert_F^2}{r_s(\widetilde U)}.
  \label{eq:qmax_stable_rank}
\end{equation}

\begin{corollary}[Conditional convergence-margin enlargement]
\label{cor:rotation_margin}
Let $q_{\max}=({\gamma}/{\mu})\lVert U\rVert_2^2$ denote the corresponding
unrotated spectral factor. For nonzero $U$ and $\widetilde U$,
\begin{equation}
  \frac{q_{\max}^{R}}{q_{\max}}
  =\frac{r_s(U)}{r_s(\widetilde U)}.
  \label{eq:qmax_ratio}
\end{equation}
Hence, if row-wise position rotation increases stable rank, then it strictly
decreases the spectral factor. Whenever the rotated iteration is contractive,
this increases its Neumann convergence margin $1-q_{\max}^{R}$; it may also
move a previously non-contractive iteration into the contractive regime.
\end{corollary}
\begin{proof}
Apply \cref{eq:qmax_stable_rank} to the rotated and unrotated bases and use
\cref{eq:frobenius_preserved}. The remaining claims follow directly from the
Neumann condition in \cref{eq:neumann_condition}.
\end{proof}

The corollary does not claim that position rotation must increase stable rank.
It isolates a precise consequence if the position-dependent rotation
redistributes a fixed amount of relation energy more evenly across singular
directions. In the inner-learning view, the nonzero singular directions of
$\widetilde U$ are the directions available to the sample-specific regression.
A larger stable rank therefore corresponds to a less concentrated effective
adaptation geometry, but not automatically to better task performance. This
conditional mechanism is testable through layer-wise stable rank, condition
number, fast-weight norm, ridge residual, and $q_{\max}^{R}$ diagnostics.

\begin{corollary}[Structure preservation under rank rotation]
\label{cor:rrlsso}
All results in
\cref{prop:spd,prop:woodbury,prop:spectrum,prop:bounds,thm:energy,thm:inner_learner,prop:inner_limit}
hold for \RRLSSO\ with $U$ replaced by $\widetilde U$. The Neumann result
additionally holds whenever
$\alpha\lambda_{\max}(\widetilde U\widetilde U^\top)<1$.
\end{corollary}
\begin{proof}
Position-dependent rotation changes the entries of the basis but not its factorized Gram form: $\widetilde U\widetilde U^\top\succeq0$. The preceding arguments apply unchanged.
\end{proof}

Together, \cref{eq:relative_kernel,cor:rrlsso} show that relative position changes the learned relation geometry without sacrificing well-posedness. The rotation is not a similarity transform of the full token-space operator, because different rows generally use different rotations; it is instead a structured reparameterization of both the Gram factors and the design matrix of the per-sample inner problem. In short, the solve removes inner optimization depth, whereas rank rotation reshapes inner adaptation geometry.
